\documentclass[11pt]{article}
\usepackage{amsmath,amssymb,amsthm,hyperref}
\newtheorem{lemma}{Lemma}
\title{Adaptive Haar residues with block biases}
\author{}
\date{}
\begin{document}
\maketitle

\paragraph{Status.}
This note records an exact reduction, a finite common-pool certificate,
and a finite Laurent-series continuation certificate. It does not prove
a compiler of unbounded superlinear capacity. All verification here is
written or exact computational verification, not Lean verification.
The original HeadComplexity definitions and admissibility conditions
are unchanged.

\paragraph{Family.}
Let $H=2^d$. Index the butterfly switches by
$e=(k,s,g)$ with $0\leq k<d$, $0\leq s<2^k$, and
$0\leq g<H/2^{k+1}$. A target word $\theta$ selects its butterfly
permutation $\pi_\theta$. Set
\[
 p_y(\epsilon)=\sum_{j=0}^{d-1}y_j\epsilon^{2^j},\qquad
 t_{k,s,g}=1-p_s,\qquad
 u_{k,g}(h)={\bf1}_{h\mathbin{\mathrm{div}}2^{k+1}=g}(-1)^{h_k}.
\]
The residues $c_y$ and the $H-1$ block biases $\beta_{k,g}$ may depend
on the entire target. The required inequalities are
\[
 (1-2\theta_{k,s,g})\left(
 \sum_h\frac{u_{k,g}(h)c_{\pi_\theta(h)}}{t_{k,s,g}-p_{\pi_\theta(h)}}
 +\beta_{k,g}\right)>0.                                      \tag{1}
\]
The query profiles and times are fixed once $d$ and $\epsilon$ are fixed.

\begin{lemma}[Exact elimination of a block bias]
Write $F_s$ for the reciprocal sum in a fixed block. There exists a real
$\beta$ giving all its required signs if and only if
$F_a>F_b$ for every $a$ with $\theta_a=0$ and $b$ with $\theta_b=1$.
If both classes occur, one may take
$\beta=-(\min_{\theta_a=0}F_a+\max_{\theta_b=1}F_b)/2$.
A block containing just one class places no restriction on the residues.
\end{lemma}
\begin{proof}
The inequalities ask for $-\min F_a<\beta<-\max F_b$ when both
classes occur. If only the positive class occurs, use
$\beta=1-\min F_a$; if only the negative class occurs, use
$\beta=-1-\max F_b$.
\end{proof}

Since all denominators are positive, each mixed-block comparison is
equivalently
\[
 \sum_y \frac{\operatorname{sgn}(p_a-p_b)
        u_{k,g}(\pi_\theta^{-1}(y))c_y}
 {(1-p_a-p_y)(1-p_b-p_y)}>0.                                  \tag{2}
\]
Indeed $F_a-F_b$ equals the displayed sum times $|p_a-p_b|$.
Thus feasibility reduces to a homogeneous rational linear system in
only $H$ variables. Rescaling a strict solution turns its finitely many
inequalities into inequalities with right side one.

\paragraph{Finite common-pool result at eight wires.}
At $d=3$ and $\epsilon=1/100$, there is a single signed residue vector
$c$ for which (1) holds for every one of the $4096$ butterfly targets,
using separate biases for the targets. The exact vector is in
\texttt{butterfly-adaptive-haar-common-h8.json}.
The solver-free checker
\texttt{butterfly-adaptive-haar-common-check.py} verifies all
$49\,152$ original strict query inequalities and finds exactly forty
distinct comparison rows (2). This is a finite-scale result. It asserts
neither arbitrary smallness of $\epsilon$ nor the corresponding
statement at larger $H$. The obstruction in
\texttt{butterfly-common-pool-bias-bound.tex} explains why a universal
compact common pool cannot be the general induction invariant.

\paragraph{Adaptive thirty-two-wire samples.}
The script \texttt{butterfly-adaptive-haar-lp.py} generates exact rational
witnesses using an ordinary continuous rational LP. It never declares
integer variables. All accepted witnesses are independently checked
using (1), with biases reconstructed by the lemma, by
\texttt{butterfly-adaptive-haar-sample-check.py}.
At $d=5$, $\epsilon=1/100$, and exponents $(1,2,4,8,16)$, ten
pseudorandom targets and three targets having a fully alternating
sixteen-query root word all have exact witnesses. These thirteen tests
provide no probability bound or universal statement. The committed
sample file stores the first four random witnesses; the first of these
also has the stronger continuation certificate below.

\paragraph{An exact Laurent continuation.}
The file \texttt{butterfly-adaptive-haar-h32-laurent.json} specifies one
of the eighty-bit targets and a square active basis. The basis consists
of comparison equalities from (2), with right side one, and selected
coordinate equalities $c_y=0$. Regard its coefficients as rational
functions of $\epsilon$. The accompanying checker
\texttt{butterfly-adaptive-haar-laurent-check.py} proves that this basis
is nonsingular over $\mathbb Q((\epsilon))$ and that every comparison
score has positive leading coefficient. Consequently this one target
satisfies (1) for every sufficiently small positive $\epsilon$.

Here are sufficient details to audit what the computation establishes.
The checker uses exact rational power series with tracked truncation
precision. It performs Gaussian elimination with full pivoting by
lowest valuation. Every pivot has a nonzero leading coefficient.
Multiplying the right side by $\epsilon^{85}$ keeps intermediate
expressions in the power-series ring. The pivot valuations are twenty-four
zeros followed by
\[
 2,2,6,6,8,8,16,22.
\]
After solving, the checker substitutes the result into the original
basis equations and verifies zero residual to the available precision.
For all $104$ comparison rows, the computed leading coefficient is
strictly positive and occurs before the tracked precision ends.
All these statements are exact rational calculations. Nonzero pivots
certify an invertible rational-function matrix; precision tracking
makes its computed leading solution terms and score terms valid in the
formal Laurent field. Each rational score with positive leading term
is positive on some punctured interval to the right of zero. Intersect
the finitely many intervals and apply the bias-elimination lemma.

The nonzero residue leading terms are listed below as triples
$(y,\text{coefficient},\text{valuation})$:
\[
\begin{gathered}
 (16,-8,-8),\ (18,1,0),\ (20,-4,-22),\ (21,4,-8),\
 (22,8,-7),\ (23,-8,-8),\ (24,-4,-8),\ (25,4,-18),\
 (26,-4,-16),\ (27,4,-16),\ (28,4,-22),\ (29,-4,-18),\
 (30,-4,-8),\ (31,-12,-8).
\end{gathered}
\]
The other residues vanish identically by their selected basis equations.
The checker also verifies this displayed leading-term data against the
certificate. These simple leading coefficients suggest a possible
recursive cancellation rule, but no such rule has been proved.

\paragraph{Remaining quantitative gap.}
A proof that (1) is feasible for all targets at arbitrarily large $H$,
or a sufficiently strong uniform lower bound on the density of feasible
targets, is absent. The construction has $H\log_2 H/2$ queries and
$H-1$ block biases, so such a result would be relevant to a compiler
with a superconstant saving. Finite samples, even with exact Laurent
continuation, do not supply that missing density or induction theorem.
\end{document}
